\section{Experiments}
\label{sec:exp}

\paragraph{Data and protocol.}
\textbf{HELM MMLU} v1.13.0 (release summary dated 2025-01-10): 79 models $\times$ 57 subjects, exact-match accuracy in points, no missing cells; the item count per subject ranges from 100 to 1534. \textbf{Open LLM Leaderboard} (\texttt{contents}, commit \texttt{9c09a7c\ldots}, 2025-03-20): 4576 rows reduced to 4497 models after removing duplicate names (best variant kept) and to 4496 after removing the one flagged model; $n=6$ normalised task scores (the leaderboard's own normalisation, which removes random-guess baselines before the official arithmetic mean). Models on both boards are not independent (merges, same-base fine-tunes); we say where this matters. The primary targets are fixed in advance as the ten challengers ranked 2--11 under the mean rule (``top-10''); results by initial rank use all targets, and every interval is a 95\% percentile bootstrap over targets unless stated. Cap scenarios and costs are those of Section~\ref{sec:method}; \emph{gain}$+h$ allows at most $h$ points per subject, \emph{front}$+c$ the best existing score on the subject plus $c$. The margin is $\varepsilon=10^{-3}$ points unless a statistical margin is stated.

\paragraph{The closed forms are exact.}
Proposition~\ref{prop:closed} agrees with the exact solvers on every target checked: for the mean, the minimum and the odd-$n$ median, the maximum absolute error in $L_1$ and $K$ is $0$ over 78--79 HELM and 4495 OLLB targets (and $\le 10^{-7}$ for $L_\infty$, the bisection tolerance); the cost of every checked target never decreases when the cap is tightened (0 violations in 1896 HELM and 3600 OLLB checks; Table~\ref{tab:theory}, appendix).

\paragraph{Rule $\times$ cap: cost differs strongly, feasibility depends on the cap.}
\begin{table}[t]
\centering\small
\caption{\protect\input{../tables/tab_helm_top10.caption.tex}}
\label{tab:helm-top10}
\resizebox{\linewidth}{!}{\input{../tables/tab_helm_top10.tex}}
\end{table}
Table~\ref{tab:helm-top10} and Figure~\ref{fig:rulecap} give the headline on HELM. Under cap \emph{ceil}, a mean-rule win costs a median of 162 points in total (CI 42--222; $K=4$) for the ten closest challengers, i.e.\ $n\cdot(\mathrm{gap}+\varepsilon)$ by Proposition~\ref{prop:closed}, which is why $L_1$ under the mean scales with the number of tasks. The median rule needs 4.7 points ($K=4.5$), the minimum rule 5.3 ($K=2$), plurality 5.2 ($K=7$), the Condorcet requirement 22.6 ($K=17$) and Borda 54.6 ($K=17.5$). The paired per-target ratio to the mean rule (Fig.~\ref{fig:rulecap}c) has median 0.052 for the median rule, 0.055 for the minimum, 0.034 for plurality, 0.17 for the Condorcet requirement, 0.41 for Borda, and 0.91 and 0.84 for the geometric and harmonic means (CIs in the figure; with ten targets they are wide, e.g.\ 0.002--0.075 for the median). $K$ and $L_1$ measure effort in score space, not training cost. Across the six caps this ordering of rules is nearly unchanged: the median ratios stay within narrow bands (median rule 0.052--0.055, minimum 0.032--0.055, Condorcet requirement 0.17--0.21, Borda 0.37--0.43), and the exception is plurality, which rises from 0.034 to 0.16 under \emph{front}$+0$. What \emph{does} depend on the cap is feasibility (Fig.~\ref{fig:rulecap}a): under \emph{gain}$+5$ first place is attainable for only 6/10 targets under the minimum rule, 8/10 under the geometric and harmonic means and 9/10 under Borda, against 10/10 under the mean (at \emph{gain}$+10$: 7/10 for the minimum). The cost medians are over feasible targets only, so they must be read together with these rates.
\begin{figure}[t]
\centering
\includegraphics[width=\linewidth]{../figures/fig_rule_cap_top10.pdf}
\caption{\protect\input{../figures/fig_rule_cap_top10.caption.tex}}
\label{fig:rulecap}
\end{figure}

\paragraph{Official HELM rule: item-count weighting.}
The official HELM ``All Subjects'' score weights subjects by their item count. For the same top-10 targets the paired $L_1$ ratio of this rule to the equal-weight mean is 0.16 (CI 0.14--0.33) under \emph{ceil}, 0.28 under \emph{gain}$+10$ and 0.47 under \emph{front}$+2$ (one-subject investment suffices under \emph{ceil}, $K=1$): a target strong on heavy subjects needs far less under the official rule, and the advantage shrinks as caps tighten.

\paragraph{Cost against initial rank.}
\begin{figure}[t]
\centering
\includegraphics[width=0.92\linewidth]{../figures/fig_rank_ceil.pdf}
\caption{\protect\input{../figures/fig_rank_ceil.caption.tex}}
\label{fig:rank}
\end{figure}
Figure~\ref{fig:rank} shows cost against the target's initial rank under the mean rule (cap \emph{ceil}). On HELM, $L_1$ increases steadily with rank for every rule, but rank rules saturate in breadth: plurality never needs more than 9 subjects and the Condorcet requirement never more than 29 (the majority of 57), while Borda needs up to 54 and the mean up to 46 (every target is feasible under \emph{ceil}). For targets ranked 31 or lower the median rule has median $K=26$, near the majority bound of Proposition~\ref{prop:closed}; the maximum over HELM targets is $K=29=(n+1)/2$. On the Open LLM Leaderboard ($n=6$, score rules only) a handful of tasks already suffices at every rank: $K\le 6$ for all rules and targets, medians 2--3 (for the ten closest challengers, mean: $K=1$, $L_1=25$ points; median rule: $L_1=0.9$; Table~\ref{tab:ollb-top10}).

\paragraph{Does breadth grow with the number of tasks?}
\begin{figure}[t]
\centering
\includegraphics[width=\linewidth]{../figures/fig_taskscale_gain10.pdf}
\caption{\protect\input{../figures/fig_taskscale_gain10.caption.tex}}
\label{fig:taskscale}
\end{figure}
We drew $n\in\{3,6,12,24,48\}$ subjects from the 57 (200 draws each) and priced the model ranked 2 or 5 \emph{under the mean rule} in that subset for every rule, so that all rules price the same models (Figure~\ref{fig:taskscale}; \emph{front}$+2$ in the appendix). For the rank-5 target under \emph{gain}$+10$ the median $K$ at $n=3,6,12,24,48$ is $1,2,4,7,13$ for the mean (slope $0.84$), $1,1,2,3,6$ for the median rule (slope $0.54$), $1,2,2,2,2$ for the minimum (slope $0.16$), $2,2,3,4,6$ for plurality ($0.51$) and $2,2,4,8,16$ for Borda ($0.82$); for the rank-2 target $K=1$ under the mean, the median and the minimum at every $n$. The total points $L_1$ grow linearly in $n$ for the mean (slope $1.00$, 0.98--1.02) and the geometric mean ($0.98$), but hardly for the median ($0.33$), the minimum ($0.06$) and plurality ($0.14$). The minimum has no feasible win in 10--31\% of the draws, and the medians are over feasible draws. Slopes are of the log mean cost over feasible draws on $\log n$ (95\% bootstrap intervals over draws: 0.82--0.86 for the mean rule and 0.51--0.56 for the median rule in $K$); slopes of the integer medians are 0.92, 0.68 and 0.20. Two conclusions follow. First, the mean rule's $K$ \emph{does} grow with $n$ once headroom per task is capped ($K\approx\mathrm{need}/h$ with $\mathrm{need}=n\cdot$gap), contrary to our prior expectation that $K$ under the mean is independent of $n$; it is independent only when one task has the headroom (cap \emph{ceil}). Second, the order-statistic rules do need more tasks as $n$ grows (median rule: yes), but far fewer than the worst case $\lfloor n/2\rfloor+1$ (6 of 48, not 25) and never more than the mean; the minimum hardly forces breadth (slope $0.16$). Our expectation that the median and the minimum scale linearly in $K$ is therefore not borne out for near-top targets: scores are strongly correlated across subjects, so such targets already clear the rival median on many subjects; the median-rule $K$ is $\approx (n+1)/2-n_{\uparrow}$, with $n_\uparrow$ the number of subjects where the target already reaches the strongest rival's median (correlation $0.95$ over non-leader targets).

\paragraph{Robustness.}
\begin{figure}[t]
\centering
\includegraphics[width=0.9\linewidth]{../figures/fig_robust.pdf}
\caption{\protect\input{../figures/fig_robust.caption.tex}}
\label{fig:robust}
\end{figure}
\emph{Noise.} We ran a parametric bootstrap of evaluation noise (100 draws of every per-subject score using its binomial standard error) and re-solved (cap \emph{gain}$+10$, top-10 targets); the median over targets of the coefficient of variation of $L_1$ across draws is 0.17 for the mean, 0.18 for the geometric mean, 0.34 for Borda and for the Condorcet requirement, 0.44 for plurality, 0.68 for the median and 0.75 for the minimum, and the minimum has no feasible win in 44\% of draws (Fig.~\ref{fig:robust}). The coefficient of variation is inflated for small costs; in points, the median 5--95\% range of $L_1$ is 89 (mean), 97 (geometric mean), 56 (Borda), 24 (Condorcet), 15 (median and minimum) and 10 (plurality), against median costs of 174, 170, 50, 22, 6.5, 7.0 and 7.6: small, but about twice the cost for the median and minimum rules. \emph{Task bootstrap} (50 draws over subjects) gives the same picture (minimum infeasible in 30\% of draws, coefficient of variation 0.58 for the median). \emph{Statistical margin.} Requiring the target to lead by $k=1$ or $2$ standard errors of the difference (0.6 points per $k$ on HELM; approximate for the median and the minimum) leaves the paired ratios nearly unchanged at \emph{ceil}: median rule 0.065 and 0.078 against 0.052 at $\varepsilon=10^{-3}$, minimum 0.055 and 0.051, geometric mean 0.91 throughout. On the Open LLM Leaderboard (margin $k\times0.96$ points) the median rule's ratio to the mean at \emph{ceil} rises from 0.033 to 0.088 and 0.127, the minimum's falls from 0.61 to 0.44 and 0.41. The margin variants were run for score rules only; for rank rules a significance margin is not defined by $\varepsilon$, and their noise behaviour is the resampling experiment above.

\paragraph{Open LLM Leaderboard: caps, empirical envelope and score normalisation.}
Cap dependence on the six-task board is stronger than on HELM (Figures~\ref{fig:rulecap-ollb} and~\ref{fig:feasible-ollb}, appendix), because six tasks leave little room to spread a gain (top-10 targets, normalised scores; \texttt{summary\_ollb\_by\_rule\_group.csv}). Under \emph{gain}$+5$ the share of targets for which first place is attainable is 1.0 for the mean and the median but 0.4 for the geometric and harmonic means, 0.3 for the minimum and 0.3 for the trimmed mean; under \emph{gain}$+10$ it is 0.7 for the minimum and, on the sampled targets below, low for plurality. The ratio to the mean-rule cost at \emph{ceil} is $0.03$ for the median, $0.23$ for Borda, $0.51$ for the Condorcet requirement, $0.61$ for the minimum, $0.64$ and $0.76$ for the harmonic and geometric means, $0.73$ for plurality and $0.88$ for the trimmed mean (under \emph{emp95}, and at \emph{ceil} on the top-500 slice); the ordering of the cheap rules is the same as on HELM, but only the median is much cheaper than the mean. \emph{Empirical envelope.} The \emph{emp95} cap is the per-task 95th percentile of the gains between 540 fine-tuned models and their bases on this leaderboard (applied to HELM it would be a transfer assumption, so we do not). Under \emph{emp95}, first place is attainable for 10/10 top-10 targets under the mean, the median, the trimmed mean and Borda, 9/10 under the geometric mean and plurality, 5/10 under the Condorcet requirement and 4/10 under the harmonic mean and the minimum; under \emph{emp99} the minimum is attainable for 5/10 and every other rule for 10/10. Even when the mean-rule win is attainable it needs $4.2$ points per task, whether that is plausible depends on the envelope assumption, not on the data. \emph{Raw versus normalised scores.} Paired per target (cap \emph{ceil}; raw minus normalised ratio $L_1^{\mathrm{rule}}/L_1^{\mathrm{mean}}$, 95\% bootstrap intervals over targets; \texttt{summary\_ollb\_raw\_vs\_norm\_paired.csv}), the median rule's ratio rises from $0.13$ to $0.84$ for ranks 11--30 ($+0.70$, 0.61--0.82) and from $0.34$ to $0.78$ for ranks 31--100 ($+0.45$, 0.43--0.46), but not for the 4396 targets beyond rank 100 ($+0.008$); for the ten closest challengers it rises from $0.03$ to $1.10$ ($+0.93$), with an interval ($-0.15$ to $1.27$) that includes zero. The minimum is dearer under raw scores in every group ($+0.24$ beyond rank 100); geometric and harmonic means have mixed signs (Figure~\ref{fig:variants}a, Table~\ref{tab:ollb-robust}, appendix). The score scale thus changes which score rule is cheap for near-top models; rank rules are invariant to per-task affine rescaling. Slices (no merged models; top 500) leave the other ratios unchanged and move the minimum (0.49), Borda (0.19) and plurality (0.51).


\paragraph{Ordinal rules on the Open LLM Leaderboard.}
Exact mixed-integer solutions for Borda, plurality and the Condorcet requirement on 4496 models are expensive, so these rules were run on a sample of 102 targets (the 20 top challengers plus 100 targets uniform in log-rank) at \emph{ceil}, \emph{front}$+2$, \emph{gain}$+10$ and \emph{emp95}. On the sample, first place is unattainable under \emph{gain}$+10$ for 37 targets under Borda, 62 under the Condorcet requirement and 89 under plurality, and under \emph{emp95} for 43, 86 and 63; among the ten closest challengers plurality is attainable for only 3 under \emph{gain}$+10$ and the Condorcet requirement for 5 under \emph{emp95}. The sample favours low-ranked models. The original board has no Condorcet winner under the six-task majority relation (HELM has one), so the Condorcet requirement there is a proxy and an upper bound on Copeland, Schulze and Kemeny.
